%%% FIELD sep-proposed
For every \(n\ge4\), the explicit model \(\mathfrak F_n\) belongs to \(\mathcal P_{\mathsf A}(1/40,1/2,32)\). Every fixed-coordinate covariance has spectrum in \([1/32,32]\); the transcript \(\mathfrak T_{\mathsf A}(P)\) has exactly \(n^2+2n-2\) matrix occurrences; every reached positive eigenvalue is at least \(1/40\); every reached graph comparison is zero or one; and both the printed source scale and the valid reciprocal scale satisfy \(\gamma^{\ast}(G)=\rho^{\ast}(G)=1\). The fixed source thresholds obey \(\gamma_{\mathrm{src}}=1/4\in(0,\gamma^{\ast}(G))\), whereas \(\eta_{\mathrm{src}}=1/80\) does not satisfy the source rank-threshold inequality because the subset \(\{3,4\}\) is not reached and
\[
0<\eta^{\ast}(P)\le\lambda_{\min}^{+}(R(\{3,4\}))\le2^{-2n-1}<1/80.
\]
Thus the family does not contradict the source theorem's existential premise that a threshold in \((0,\eta^{\ast}(P))\) may be chosen. Instead, it exhibits constant reached gaps together with a positive all-subset margin that tends to zero exponentially with \(n\). The family also satisfies every non-margin condition of Acartürk et al.'s graph theorem: the induced support-coordinate densities and causal conditionals are strictly positive and continuously differentiable, there is one distinct single-target intervention per node, their Assumption 1 parent-ratio condition holds, and the empirical Gaussian covariance score-difference estimator satisfies their consistency condition (15).
%%% FIELD strict-proposed
Restrict attention to the linear Gaussian perfect variance-reset subclass satisfying the non-margin hypotheses used by Acartürk et al., the observed-conditioning bound \(\kappa\), and independent environment sampling. If equation (22) obeys \(\eta^{\ast}(P)\ge g_0\) and the actual reached comparisons obey \(c_{t,j}(P)\in\{0\}\cup[h_0,1]\), then \(P\in\mathcal P_{\mathsf A}(g_0,h_0,\kappa)\), so Theorem 1 supplies the stated Gaussian recover-or-abstain certificate. This implication is confined to that Gaussian covariance subclass and does not transfer the positive certificate to the full smooth soft-intervention class of the source. Conversely, for every \(n\ge4\), \(\mathfrak F_n\) is both a legal member of that published class satisfying its Assumption 1 and other stated non-margin conditions and a member of \(\mathcal P_{\mathsf A}(1/40,1/2,32)\), while \(0<\eta^{\ast}(P)\le2^{-2n-1}<1/80\). Consequently, for every constant \(c>0\), some legal member of the published class is recoverable by the constant-reached-gap Gaussian certificate and has \(0<\eta^{\ast}(P)<c\). Hence no positive dimension-uniform constant lower bound on equation (22)'s all-subset margin is necessary for recovery. The source theorem's existential threshold premise remains satisfied, but on this family its equation-(22)-controlled sufficient analysis must resolve an exponentially small scale, whereas the reached-gap certificate retains constant spectral and projector gaps; its positive inverse-power sample bound is therefore exponentially pessimistic on this family. The conclusion uses actual projector geometry: the printed equation (23) quantity \(\gamma^{\ast}(G)=\min_i\|H_{i,:}\|_2\|G_{:,i}\|_2\) is not a projector-product norm, whereas Proposition 1 gives the valid reciprocal witness bound \(\rho_i(G)\).
%%% FIELD sep-proof
Put \(\epsilon=2^{-n}\).  The structural matrix of \(\mathfrak F_n\) has nonzero rows only at labels \(3\) and \(4\), with parent vectors \(b_3=(1,1)\) and \(b_4=(1,1+\epsilon)\).  Thus \((B^{(m)})^2=0\) and \(\|B^{(m)}\|_{\mathrm{op}}\le\|B^{(m)}\|_{\mathrm F}<3\) in every environment, because an intervention only deletes one row.  By the structural equations and the perfect-intervention construction,
\[
\Sigma_m=(I-B^{(m)})^{-1}\Omega_m(I-B^{(m)})^{-\top},
\qquad
(I-B^{(m)})^{-1}=I+B^{(m)},
\]
and \(I\preceq\Omega_m\preceq2I\).  Therefore
\[
 \lambda_{\max}(\Sigma_m)
 \le2\|I+B^{(m)}\|_{\mathrm{op}}^2<32,
 \qquad
 \lambda_{\min}(\Sigma_m)
 \ge\|(I-B^{(m)})^{\top}\|_{\mathrm{op}}^{-2}>1/16.
 \tag{20}
\]
In particular every fixed-coordinate covariance has spectrum in \([1/32,32]\).  The construction directly verifies the linear Gaussian DAG equations, graph compatibility, one perfect variance-reset intervention at each target, and full-rank mixing with \(G=I_n\).  Independent streams with these environment laws satisfy the sampling condition.

We next compute the primitive matrices from their definition.  For an intervention at a root \(i\in\{1,2,5,\ldots,n\}\), only the corresponding noise precision changes from \(1\) to \(1/2\).  Hence \(\Delta_i=e_ie_i^{\top}/2\), and because \(\mathbb E[Z_i^2]=1\),
\[
 R_i=\Delta_i\Sigma_0\Delta_i=\frac14e_ie_i^{\top}.
 \tag{21}
\]
For a child \(c\in\{3,4\}\) with parent vector \(b\), write \(s=b^{\top}(Z_1,Z_2)\) and \(E=Z_c-s\).  Under the observational law, \(s\) and \(E\) are independent, \(\mathbb E[E^2]=1\), and \(\mathbb E[s^2]=\|b\|_2^2\).  Direct subtraction of the observational and intervened Gaussian precision matrices gives
\[
d_c(Z)=(-bE,(E-s)/2)
\]
on the parent coordinates and coordinate \(c\).  Taking the observational second moment therefore yields
\[
 R_c=
 \begin{pmatrix}
 bb^{\top}&-b/2\\
 -b^{\top}/2&(1+\|b\|_2^2)/4
 \end{pmatrix},
 \qquad
 \operatorname{col}(R_c)=\operatorname{span}\{b,e_c\}.
 \tag{22}
\]
In the orthonormal basis \(b/\|b\|_2,e_c\), put \(B_0=\|b\|_2^2\in[2,5]\).  The nonzero block in (22) has determinant \(B_0^2/4\) and trace \((5B_0+1)/4\).  Since the smaller eigenvalue of a positive-definite \(2\)-by-\(2\) matrix is at least its determinant divided by its trace,
\[
 \lambda_{min}^{+}(R_c)
 \ge\frac{B_0^2}{5B_0+1}\ge\frac4{11},
 \qquad
 R_c\succeq\frac4{11}P_{\operatorname{span}\{b,e_c\}}.
 \tag{23}
\]

We now bound every subset sum except the exceptional child pair.  Root-only sets have positive eigenvalues \(1/4\), and a child singleton has positive eigenvalues at least \(4/11\) by (23).  If a subset contains a root \(i\in\{1,2\}\) and a child with parent vector \(b\), then on the parent plane
\[
 \lambda_{min}^{+}(P_{e_i}+P_b)
 =1-\frac{|e_i^{\top}b|}{\|b\|_2}
 \ge1-\frac2{\sqrt5}>\frac1{10}.
\]
Combining (21) and (23) shows that the parent-plane eigenvalues are greater than \(1/40\), while (23) supplies the orthogonal child coordinate with coefficient at least \(4/11\).  Adding the other child cannot lower any existing Rayleigh quotient and supplies its new orthogonal child coordinate; adding the other root supplies \((1/4)I_2\) on the parent plane.  These cases exhaust every nonempty subset of \(\{1,2,3,4\}\) other than \(\{3,4\}\).  Each label at least \(5\) contributes the orthogonal summand \(e_je_j^{\top}/4\).  Consequently every positive eigenvalue of every such subset sum is at least \(1/40\).

The population leaf-removal order under the smallest-label rule is
\[
 (3,4,1,2,5,6,\ldots,n),
\]
and hence the reversed topological order is \(\pi=(n,n-1,\ldots,5,2,1,4,3)\).  Indeed, in the first order round, deleting either child lowers the range dimension by one, whereas deleting either root leaves both linearly independent parent vectors \(b_3,b_4\) and does not lower it.  The smallest admissible child is \(3\).  After label \(3\) is removed, the same argument selects \(4\); thereafter every remaining primitive has a distinct coordinate range and the smallest-label rule gives \(1,2,5,\ldots,n\).

In the first order round, deleting one label retains a core root whenever both children remain; after label \(3\) is removed, no later order query contains both children.  In any graph query containing both children, both core roots already lie in the relevant prefix, descendant deletion removes no root, and deletion of the current row removes at most one root.  A singleton encoder query contains only one child.  Thus no reached matrix occurrence has core intersection exactly \(\{3,4\}\).  The preceding subset bounds prove that every reached positive eigenvalue is at least \(1/40\).

For each graph query, descendant pruning leaves an ancestral set.  Its primitive ranges include the coordinate of each retained root; whenever a retained child has a parent direction not already supplied by a root, (22) supplies that direction together with the child's own coordinate.  Removing the current label deletes its coordinate exactly in the zero-comparison case and otherwise leaves it supplied by a retained descendant.  Hence the two ranges appearing in each graph comparison are coordinate subspaces, and their projector-product norm is respectively zero or one according as the relevant inclusion holds.  This proves the reached projector band with \(h_0=1/2\), using the actual orthogonal projectors; it does not use the source's printed angular expression.

The order phase evaluates \(\sum_{t=2}^nt=n(n+1)/2-1\) matrix occurrences.  The graph phase contributes \(n-1\) prefix occurrences and \(n(n-1)/2\) comparison-sum occurrences, and the encoder phase contributes \(n\) singleton occurrences.  Therefore the total is
\[
 n(n+1)/2-1+(n-1)+n(n-1)/2+n=n^2+2n-2.
 \tag{24}
\]

For the unqueried subset, the four range vectors \(e_3,b_3,e_4,b_4\) are linearly independent because \(e_3,e_4\) are orthogonal to the parent plane and \(\det(b_3,b_4)=\epsilon\ne0\).  Thus \(R(\{3,4\})\) is positive definite on the four-dimensional core.  With \(v=(1,-1,0,0)^{\top}/\sqrt2\), equation (22) gives
\[
 v^{\top}R_3v=0,
 \qquad
 v^{\top}R_4v=\epsilon^2/2.
\]
The Rayleigh variational principle therefore yields
\[
 0<\eta^{\ast}(P)
 \le\lambda_{min}^{+}(R(\{3,4\}))
 \le\epsilon^2/2=2^{-2n-1}<1/80.
 \tag{25}
\]
The strict positivity on the left also follows directly because there are finitely many subset sums and every nonzero positive-semidefinite subset sum has finitely many strictly positive eigenvalues; the preceding case analysis covers all of them.  Thus a threshold smaller than \(\eta^{\ast}(P)\) exists for each fixed \(n\), even though no fixed positive threshold works uniformly in \(n\).  Since \(G=H=I_n\), the source norm-product scale and the valid reciprocal scale both equal one: \(\gamma^{\ast}(G)=\rho^{\ast}(G)=1\).  Hence \(\gamma_{\mathrm{src}}=1/4\) lies below both scales, whereas the fixed choice \(\eta_{\mathrm{src}}=1/80\) lies above the margin in (25).

Finally, every induced Gaussian density and causal conditional is strictly positive and continuously differentiable.  Let \(\phi_v\) denote the centered Gaussian density of variance \(v\).  A root intervention has a nonconstant ratio of centered Gaussian densities.  For a child, the intervention-to-observational conditional-density ratio is
\[
 \frac{\phi_2(z_c)}{\phi_1(z_c-b^{\top}z_{\operatorname{pa}(c)})}.
\]
Its derivative in either parent coordinate is not the zero function because both entries of \(b_3\) and \(b_4\) are nonzero.  This verifies the source's Assumption 1 parent-ratio condition, and a perfect stochastic intervention is a legal special case of its smooth soft-intervention model.  Entrywise laws of large numbers give \(S_{m,r}\to\Sigma_m\) in probability.  On the uniformly positive-definite neighborhood established in (20),
\[
 S_{m,r}^{-1}-\Sigma_m^{-1}
 =S_{m,r}^{-1}(\Sigma_m-S_{m,r})\Sigma_m^{-1}
\]
gives inverse convergence, and multiplication by a Gaussian vector with finite second moment gives the source's squared-mean score-difference consistency.  Thus all promised-class assumptions and all stated non-margin source conditions hold.  Equations (20)--(25) prove the claimed constants and also show precisely why the example leaves the source's existential threshold premise intact.
%%% FIELD strict-proof
On the stated Gaussian variance-reset subclass, equation (22), as recorded in \(\mathrm{def{:}source\mbox{-}margins}\), gives
\[
 \lambda_i(R(M))\in\{0\}\cup[g_0,\infty)
 \quad\text{for every }M\subseteq[n]\text{ and every }i\in[n].
 \tag{26}
\]
In particular (26) holds for every matrix occurrence reached by \(\mathsf A\), so it is the reached eigen band.  The separately stated condition \(c_{t,j}(P)\in\{0\}\cup[h_0,1]\) is exactly the reached projector band.  The Gaussian structural, intervention, mixing, conditioning, and sampling conditions in the theorem are the remaining member properties of \(\mathrm{def{:}promised\mbox{-}class}\).  Hence
\[
 P\in\mathcal P_{\mathsf A}(g_0,h_0,\kappa).
\]
Theorem \(\mathrm{thm{:}promised\mbox{-}transcript\mbox{-}certificate}\) now supplies the recover-or-abstain guarantee, the exact environment-labelled transitive closure on return, and the stated encoder error bar.  This implication uses the Gaussian covariance model and therefore proves no upward transfer of the positive certificate to the source's larger smooth soft-intervention class.

For the family comparison, Theorem \(\mathrm{thm{:}separation\mbox{-}family}\) proves for every \(n\ge4\) that
\[
 \mathfrak F_n\in\mathcal P_{\mathsf A}(1/40,1/2,32),
 \qquad
 0<\eta^{\ast}(P)\le2^{-2n-1}.
 \tag{27}
\]
It also proves directly that \(\mathfrak F_n\) has smooth positive causal conditionals, one legal distinct single-target intervention per node, the parent-ratio property in Assumption 1, and the required score consistency.  It is therefore a member of the published model class, not merely of a newly constructed analogue.  Theorem \(\mathrm{thm{:}promised\mbox{-}transcript\mbox{-}certificate}\) recovers on this family under constant public gaps \(g_0=1/40\) and \(h_0=1/2\).

Let \(c>0\) be arbitrary.  Because \(2^{-2n-1}\to0\), choose an integer \(n_c\ge4\) such that \(2^{-2n_c-1}<c\).  Applying (27) at \(n_c\) gives a legal, recoverable published-class member with
\[
 0<\eta^{\ast}(P)<c.
\]
Since this holds for every \(c>0\), no positive constant independent of dimension can be a necessary lower bound on the all-subset margin for recovery.  This is the exact quantifier justified by the family.  It does not say that the margin vanishes at a fixed \(n\), and it does not refute the source theorem's existential premise: by (27), the interval \((0,\eta^{\ast}(P))\) is nonempty for every fixed member.

Nevertheless, any threshold satisfying that premise on \(\mathfrak F_n\) must obey
\[
 0<\eta<\eta^{\ast}(P)\le2^{-2n-1},
 \qquad\text{so}\qquad
 \eta^{-1}>2^{2n+1}.
 \tag{28}
\]
Thus the equation-(22)-controlled sufficient analysis operates at an exponentially small resolution scale on this family.  Its displayed sample requirement has positive inverse-power dependence on that threshold, so (28) makes that sufficient sample bound exponential in \(n\), while the present Gaussian certificate uses the constant reached gaps \(1/40\) and \(1/2\).  This proves exponential pessimism of the sufficient bound on the family without asserting failure of its existential hypothesis.

The projector qualification is algebraic.  Since \(H=G^{\dagger}\), biorthogonality gives \(H_{i,:}G_{:,i}=1\), and Cauchy--Schwarz gives
\[
 \|H_{i,:}\|_2\|G_{:,i}\|_2\ge1.
\]
Every product of orthogonal projectors has operator norm at most one, so the norm product cannot be the asserted general projector-product value when it is strictly larger than one.  Proposition \(\mathrm{prop{:}projector\mbox{-}geometry}\) instead proves the valid reciprocal lower bound \(\rho_i(G)\) for any witnessed nonzero comparison.  No step above uses the false printed identity.
